\section{Thesis Outline}
\label{sec:thesis-outline}

This thesis is organized into eight chapters as follows.

\textbf{Chapter 1 -- Introduction} presents the motivation for developing Stockrium and the problem addressed by the thesis. It defines the project's primary objectives, application scope, and overall thesis structure.

\textbf{Chapter 2 -- Related Work and Foundation Technologies} reviews studies on financial decision support, large language models, multi-agent systems, news sentiment, and Vietnamese stock-market prediction. It then establishes technical, fundamental, and financial-news analysis foundations, introduces multi-agent coordination concepts together with backtesting principles, and presents the development-technology foundations required to understand Stockrium.

\textbf{Chapter 3 -- Proposed Application} describes the application idea, target stakeholders, principal use cases, and user flows. It presents the investor-facing mobile experience and the Stockrium Lab workflow for running, retaining, reviewing, and comparing supported AI-analysis and backtesting experiments.

\textbf{Chapter 4 -- System Design and Implementation} explains the overall architecture and the design of the mobile application, Stockrium Lab, backend services, storage, and data-acquisition pipelines. It also details the multi-agent stock-analysis workflow, experiment persistence and comparison, and the historical backtesting framework.

\textbf{Chapter 5 -- System Testing} defines the testing objectives, environment, and verification strategy for the implemented system. It reports functional and integration testing of the mobile and web interfaces, backend services, authentication mechanisms, multi-agent workflows, experiment management, and backtesting flow.

\textbf{Chapter 6 -- Representative Use of Stockrium Lab} follows a historical VN30 case in which a user configures a run, monitors its execution, reopens the saved result, and compares selected outputs with explicit references. The chapter retains only the values needed to demonstrate the tool and summarizes the principal interpretation limits without treating the case as proof of superior returns.

\textbf{Chapter 7 -- Demo} presents the implemented mobile and Stockrium Lab interfaces corresponding to the principal use cases, including market discovery, evidence inspection, configurable multi-agent analysis, the conversational financial assistant, experiment history and comparison, and historical backtesting.

\textbf{Chapter 8 -- Conclusion} summarizes the project outcomes, principal contributions, and degree to which the stated objectives were achieved. It discusses the remaining limitations and proposes directions for extending the data, analytical methods, experimentation capabilities, and software platform.
